% ==================================================================
%  preambolo.tex : shared preamble for esercizi.tex and soluzioni.tex
%  Italian first, English after (grey italic). Decimal comma in math.
% ==================================================================
\documentclass[11pt,a4paper]{article}

\usepackage[T1]{fontenc}
\usepackage[utf8]{inputenc}
\usepackage{amsmath,amssymb}
\usepackage{libertinus}
\usepackage{libertinust1math}
\usepackage[english,italian]{babel}
\usepackage{icomma}            % 2,5 is a decimal number in math mode
\usepackage{geometry}
\geometry{a4paper, left=2cm, right=2cm, top=2.2cm, bottom=2cm, headheight=14pt}
\usepackage{microtype}
\usepackage[table]{xcolor}
\usepackage{tikz}
\usetikzlibrary{arrows.meta, calc, angles, quotes, patterns, decorations.markings}
\usepackage[most]{tcolorbox}
\usepackage[inline]{enumitem}
\usepackage{comment}
\usepackage{multicol}
\usepackage{needspace}
\usepackage{array, booktabs, tabularx}
\usepackage{cancel}
\usepackage{fancyhdr}
\usepackage[hidelinks]{hyperref}

\setlength{\parindent}{0pt}
\setlength{\parskip}{3pt plus 1pt}
\setlist{nosep, leftmargin=1.6em, topsep=2pt, itemsep=2pt}
\setlist[enumerate,1]{label=\textbf{\alph*)}}
\newlist{elenco}{enumerate*}{1}
\setlist[elenco]{label=\textbf{\alph*)}, itemjoin=\hspace{2.2em}, afterlabel=\ }

% ---------------------------------------------------------- colours
\definecolor{navy}{HTML}{1B2A4A}
\definecolor{engray}{HTML}{55606E}
\definecolor{lv1}{HTML}{2E8B57}
\definecolor{lv2}{HTML}{D68910}
\definecolor{lv3}{HTML}{C0392B}
\definecolor{solbg}{HTML}{F1F8F3}
\definecolor{solfr}{HTML}{2E8B57}
\definecolor{trapbg}{HTML}{FDF1F0}
\definecolor{fgA}{HTML}{C0392B}
\definecolor{fgB}{HTML}{1F6FB2}
\definecolor{fgC}{HTML}{2E8B57}
\definecolor{fgD}{HTML}{D68910}
\definecolor{lightfill}{HTML}{DCE6F2}

% ---------------------------------------------------------- language
% \en{...}: the English version, on its own line, grey italic
\newcommand{\en}[1]{\par{\small\itshape\color{engray}#1\par}}
% \enx{...}: short English gloss inside a sentence
\newcommand{\enx}[1]{{\small\itshape\color{engray}(#1)}}

% ---------------------------------------------------------- math helpers
\newcommand{\ang}[3]{#1\widehat{#2}#3}
\newcommand{\av}[1]{\widehat{#1}}
\newcommand{\seg}[1]{\overline{#1}}
\newcommand{\tri}[1]{\triangle #1}
\newcommand{\unit}[1]{\,\mathrm{#1}}
\newcommand{\cm}{\unit{cm}}
\newcommand{\cmq}{\unit{cm^2}}
\newcommand{\cmc}{\unit{cm^3}}
\newcommand{\mq}{\unit{m^2}}

% ---------------------------------------------------------- levels
\newcommand{\stella}[1]{\textcolor{#1}{$\bigstar$}}
\newcommand{\livello}[1]{%
  \ifcase#1\or\stella{lv1}\or\stella{lv2}\stella{lv2}\or\stella{lv3}\stella{lv3}\stella{lv3}\fi}

% ---------------------------------------------------------- sections
\newcounter{sezione}
% \sezione{Titolo italiano}{English title}
\newcommand{\sezione}[2]{%
  \par\addvspace{14pt}\Needspace{10\baselineskip}%
  \refstepcounter{sezione}%
  \noindent\begin{tcolorbox}[enhanced, colback=navy, colframe=navy, sharp corners,
      left=6pt, right=6pt, top=3pt, bottom=3pt, boxsep=0pt, before skip=0pt, after skip=6pt]
    {\color{white}\large\bfseries \thesezione.\ #1\hfill{\normalsize\itshape #2}}
  \end{tcolorbox}\par}

% ---------------------------------------------------------- exercises
% \begin{esercizio}{ID}{level 1..3}{Titolo}{Title}  statement  \end{esercizio}
% the statement is kept in one piece (never split across pages); the solution follows it
\newtcolorbox{esercizio}[4]{blank, unbreakable, before skip=10pt plus 2pt, after skip=2pt,
  before upper={\noindent{\large\bfseries\color{navy}#1}\hspace{0.6em}\livello{#2}\hspace{0.6em}%
  {\bfseries #3}\hspace{0.4em}{\color{engray}\itshape\,/\ #4}\par
  \vspace{-3pt}{\color{navy!30}\rule{\linewidth}{0.6pt}}\par\vspace{1pt}}}

% solution box: shown in soluzioni.tex; esercizi.tex defines \modoesercizi and hides it
\ifdefined\modoesercizi
  \excludecomment{solbox}
\else
\newtcolorbox{solbox}{enhanced, breakable, colback=solbg, colframe=solfr, boxrule=0.6pt,
  arc=2pt, left=6pt, right=6pt, top=4pt, bottom=4pt, before skip=6pt, after skip=4pt,
  title={Soluzione \textperiodcentered\ \itshape Solution}, fonttitle=\sffamily\bfseries\small,
  coltitle=white, colbacktitle=solfr, attach boxed title to top left={xshift=6pt, yshift=-2pt},
  boxed title style={sharp corners, boxrule=0pt}}
\fi
% warning inside a solution: common mistake
\newtcolorbox{attenzione}{enhanced, colback=trapbg, colframe=fgA!70, boxrule=0.5pt, arc=2pt,
  left=5pt, right=5pt, top=2pt, bottom=2pt, before skip=4pt, after skip=2pt, fontupper=\small}
\newcommand{\errore}[2]{\begin{attenzione}\textbf{\color{fgA}Attenzione:} #1\par
  {\itshape\color{engray}Watch out: #2}\end{attenzione}}

% ---------------------------------------------------------- tikz styles
\tikzset{
  pt/.style={circle, fill=black, inner sep=1.1pt},
  side/.style={line width=0.9pt, line join=round},
  helper/.style={line width=0.5pt, dashed, black!55},
  hidden/.style={line width=0.6pt, dashed, black!50},
  axis/.style={-{Stealth[length=2mm]}, line width=0.7pt},
  gridl/.style={black!12, line width=0.3pt},
  tick1/.style={postaction={decorate, decoration={markings,
      mark=at position 0.5 with {\draw[line width=0.7pt] (0,-3pt)--(0,3pt);}}}},
  tick2/.style={postaction={decorate, decoration={markings,
      mark=at position 0.5 with {\draw[line width=0.7pt] (-1.2pt,-3pt)--(-1.2pt,3pt) (1.2pt,-3pt)--(1.2pt,3pt);}}}},
}
\newcommand{\rang}[3]{\pic[draw, line width=0.5pt, angle radius=2.2mm] {right angle=#1--#2--#3};}

% ---------------------------------------------------------- header
\pagestyle{fancy}
\fancyhf{}
\renewcommand{\headrulewidth}{0.4pt}
\fancyfoot[C]{\small\thepage}
